% EGE 2027, task 8. Source IDs and URLs are kept in comments.
% Editable task bank. Append new BankTask blocks at the end of the document.
\documentclass[12pt,a4paper]{article}
\usepackage{physics-textbook}
\usepackage{graphicx,needspace}
\newif\ifShowSolutions
\ifdefined\TutorEdition
  \ShowSolutionstrue
\else
  \ShowSolutionsfalse
\fi
\newcounter{banktask}
\newcommand{\BankTask}[3]{%
  \refstepcounter{banktask}%
  \Needspace{35mm}\par\medskip\noindent\begin{minipage}{\linewidth}%
  \noindent\textbf{\thebanktask.}\quad #2\par
  \ifShowSolutions
    #3\par
    \ifdefined\TutorEdition
      \ifcsname TutorSource#1\endcsname
        \smallskip{\small\csname TutorSource#1\endcsname\par}%
      \fi
    \fi
  \fi
  \end{minipage}\par\medskip
}
\begin{document}
Справочные величины (если в условии не оговорено иное):\par
$p_{\text{атм}}=10^5\,\text{Па}$\par
$c_{\text{свинца}}=130\,\text{Дж}/(\text{кг}\cdot\text{К})$\par
$\lambda_{\text{свинца}}=2{,}5\cdot10^4\,\text{Дж}/\text{кг}$



% Source BA6eFF: https://ege.fipi.ru/bank/questions.php?search=1&proj=BA1F39653304A5B041B656915DC36B38&qid=BA6eFF
\BankTask{BA6eFF}{
Двум моль газа сообщили количество теплоты, равное \(1700\ \text{Дж}\), при этом внешние силы совершили над газом работу, равную \(500\ \text{Дж}\). На сколько увеличилась внутренняя энергия газа (в $\text{Дж}$)?
}{
\smallskip\noindent\textbf{Решение от GPT}\par
По первому закону термодинамики
\[
\Delta U=Q+A_{\text{внеш}}=1700+500=2200\ \text{Дж}.
\]
Ответ: \(2200\ \text{Дж}\).
}

% Source 14FBF5: https://ege.fipi.ru/bank/questions.php?search=1&proj=BA1F39653304A5B041B656915DC36B38&qid=14FBF5
\BankTask{14FBF5}{
Газ получил количество теплоты, равное \(300\ \text{Дж}\), при этом внутренняя энергия газа уменьшилась на \(100\ \text{Дж}\). Масса газа не менялась. Какую работу совершил газ в этом процессе (в $\text{Дж}$)?
}{
\smallskip\noindent\textbf{Решение от GPT}\par
По первому закону термодинамики \(Q=\Delta U+A\), где \(A\) — работа газа. Так как \(\Delta U=-100\ \text{Дж}\),
\[
A=Q-\Delta U=300-(-100)=400\ \text{Дж}.
\]
Ответ: \(400\ \text{Дж}\).
}

% Source 5AD1F4: https://ege.fipi.ru/bank/questions.php?search=1&proj=BA1F39653304A5B041B656915DC36B38&qid=5AD1F4
\BankTask{5AD1F4}{
Газ в сосуде сжали, совершив работу, равную \(500\,\text{Дж}\). Внутренняя энергия газа при этом увеличилась на \(350\,\text{Дж}\). Какое количество теплоты отдал газ окружающей среде (в $\text{Дж}$)?
}{
\smallskip\noindent\textbf{Решение от GPT}\par
По первому закону термодинамики \(\Delta U=Q+A_{\text{внеш}}\), где \(Q\) — теплота, полученная газом. Тогда \(Q=350-500=-150\,\text{Дж}\). Газ отдал окружающей среде \(Q_{\text{отд}}=150\,\text{Дж}\).

Ответ: \(150\,\text{Дж}\).
}

% Source DE8578: https://ege.fipi.ru/bank/questions.php?search=1&proj=BA1F39653304A5B041B656915DC36B38&qid=DE8578
\BankTask{DE8578}{
При изобарном расширении постоянной массы неона при нормальном атмосферном давлении ему сообщили количество теплоты, равное \(200\,\text{кДж}\). Определите изменение объёма неона. (в $\text{м}^3$).
}{
\smallskip\noindent\textbf{Решение от GPT}\par
Неон — одноатомный идеальный газ, поэтому \(\Delta U=\frac32p\Delta V\). По первому закону термодинамики \(Q=\Delta U+p\Delta V=\frac52p\Delta V\). Принимая нормальное атмосферное давление \(p=1{,}0\cdot10^5\,\text{Па}\), получаем
\[
\Delta V=\frac{2Q}{5p}=\frac{2\cdot200\cdot10^3}{5\cdot10^5}=0{,}8\,\text{м}^3.
\]
Ответ: \(0{,}8\,\text{м}^3\).
}

% Source F23045: https://ege.fipi.ru/bank/questions.php?search=1&proj=BA1F39653304A5B041B656915DC36B38&qid=F23045
\BankTask{F23045}{
\begin{center}
\begin{tikzpicture}[task graph,x=0.65cm,y=0.65cm]
\draw[task grid,step=1] (0,0) grid (5,5);
\draw[task axis] (0,0)--(6,0) node[below] {$V$};
\draw[task axis] (0,0)--(0,6) node[left] {$p$};
\node[below left] at (0,0) {$0$};
\draw[task curve] (1,4)--(3,4)--(5,2);
\node[above left] at (1,4) {$1$};
\node[above right] at (3,4) {$2$};
\node[right] at (5,2) {$3$};
\end{tikzpicture}
\end{center}
На рисунке показано, как меняется давление идеального газа в зависимости от его объёма при переходе из состояния $1$ в состояние $2$, а затем в состояние $3$. Чему равно отношение работ газа $\dfrac{A_{23}}{A_{12}}$ в этих двух процессах?
}{
\smallskip\noindent\textbf{Решение от GPT}\par
Работа газа равна площади под графиком $p(V)$. Обозначим цену клетки по осям через $\Delta p$ и $\Delta V$. Тогда
\[
A_{12}=4\Delta p\cdot2\Delta V=8\Delta p\Delta V,
\qquad
A_{23}=\frac{4\Delta p+2\Delta p}{2}\cdot2\Delta V=6\Delta p\Delta V.
\]
Следовательно, $\dfrac{A_{23}}{A_{12}}=\dfrac{6}{8}=0{,}75$.

Ответ: $0{,}75$.
}

% Source F45AFF: https://ege.fipi.ru/bank/questions.php?search=1&proj=BA1F39653304A5B041B656915DC36B38&qid=F45AFF
\BankTask{F45AFF}{
\begin{center}
\begin{tikzpicture}[task graph,x=1cm,y=1cm]
\draw[task axis] (0,0)--(4.5,0) node[below] {$Q,\,\text{кДж}$};
\draw (0,0)--(0,0.44);
\draw[task axis] (0,0.68)--(0,3.35) node[left] {$T,\,\text{К}$};
\draw (-0.09,0.46)--(0.02,0.52)--(0.12,0.48);
\draw (-0.09,0.57)--(0.02,0.63)--(0.12,0.59);
\node[below left] at (0,0) {$0$};
\draw (0.42,-0.06)--(0.48,0.06);\draw (0.54,-0.06)--(0.60,0.06);
\draw[task projection] (1.45,0)--(1.45,1.05);
\draw[task projection] (3.1,0)--(3.1,2.6);
\node[below] at (1.45,0) {$100$};
\node[below] at (3.1,0) {$136$};
\draw[task projection] (0,1.05)--(1.45,1.05);
\draw[task projection] (0,2.6)--(3.1,2.6);
\node[left] at (0,1.05) {$273$};
\node[left] at (0,2.6) {$293$};
\draw[task curve] (0.7,0.3454545)--(3.1,2.6);
\fill (1.45,1.05) circle (1.7pt);
\fill (3.1,2.6) circle (1.7pt);
\end{tikzpicture}
\end{center}
На рисунке показана зависимость температуры металлической детали массой $2\,\text{кг}$ от переданного ей количества теплоты. Чему равна удельная теплоёмкость металла (в $\text{Дж}/(\text{кг}\cdot\text{К})$)?
}{
\smallskip\noindent\textbf{Решение от GPT}\par
Используем $\Delta Q=mc\Delta T$. По графику $\Delta Q=(136-100)\,\text{кДж}=36000\,\text{Дж}$, $\Delta T=293-273=20\,\text{К}$. Тогда
\[
c=\frac{\Delta Q}{m\Delta T}=\frac{36000}{2\cdot20}=900\,\frac{\text{Дж}}{\text{кг}\cdot\text{К}}.
\]
Ответ: $900$.
}

% Source 217002: https://ege.fipi.ru/bank/questions.php?search=1&proj=BA1F39653304A5B041B656915DC36B38&qid=217002
\BankTask{217002}{
\begin{center}
\begin{tikzpicture}[task graph,x=1.5cm,y=1cm]
\draw[task axis] (0,0)--(3.8,0) node[below] {$Q,\,10^4\,\text{Дж}$};
\draw[task axis] (0,0)--(0,3) node[left] {$t,\,{}^\circ\text{C}$};
\node[below left] at (0,0) {$0$};
\foreach \x in {0.5,1,1.5,2} {
\draw[task projection] (\x,0)--(\x,1.75);
}
\draw[task projection] (2.55,0)--(2.55,2.50625);
\draw[task curve] (0,0.8)--(0.5,1.75)--(2,1.75)--(2.8,2.85);
\node[below] at (1,0) {$1$};
\node[below] at (2,0) {$2$};
\end{tikzpicture}
\end{center}
На рисунке показан график изменения температуры вещества по мере поглощения им количества теплоты. Вещество находится в сосуде под поршнем. Масса вещества равна $0{,}5\,\text{кг}$. Первоначально вещество было в жидком состоянии. Какова удельная теплота парообразования вещества (в $\text{кДж}/\text{кг}$)?
}{
\smallskip\noindent\textbf{Решение от GPT}\par
Горизонтальный участок соответствует парообразованию. На него затрачено
\[
Q_{\text{пар}}=(2-0{,}5)\cdot10^4=1{,}5\cdot10^4\,\text{Дж}.
\]
Удельная теплота парообразования
\[
L=\frac{Q_{\text{пар}}}{m}=\frac{1{,}5\cdot10^4}{0{,}5}=3\cdot10^4\,\text{Дж}/\text{кг}=30\,\text{кДж}/\text{кг}.
\]
Ответ: $30$.
}

% Source 5745AB: https://ege.fipi.ru/bank/questions.php?search=1&proj=BA1F39653304A5B041B656915DC36B38&qid=5745AB
\BankTask{5745AB}{
\begin{center}
\begin{tikzpicture}[task graph,x=0.19cm,y=0.95cm]
\draw[task axis] (0,0)--(31,0) node[right] {$t,\,\text{мин.}$};
\draw[task axis] (0,0)--(0,4) node[left] {$T$};
\node[below left] at (0,0) {$0$};
\foreach \x in {5,10,15,20,25}{
  \node[below] at (\x,0) {$\x$};
}
\foreach \x/\y in {5/2.35,10/2.35,15/2.35,20/2.9,25/3.45}
  \draw[task projection] (\x,0)--(\x,\y);
\draw[task curve] (0,1.45)--(5,2.35)--(15,2.35)--(25,3.45);
\end{tikzpicture}
\end{center}
Для плавления куска льда при температуре его плавления требуется количество теплоты, равное \(3\,\text{кДж}\). Этот кусок льда внесли в тёплое помещение. Зависимость температуры льда от времени представлена на рисунке. Определите среднюю тепловую мощность, подводимую к куску льда в процессе плавления. (в $\text{Вт}$).
}{
\smallskip\noindent\textbf{Решение от GPT}\par
Плавлению соответствует горизонтальный участок графика от \(5\) до \(15\,\text{мин}\). Поэтому \(\Delta t=10\,\text{мин}=600\,\text{с}\), и средняя мощность равна
\[
P=\frac{Q}{\Delta t}=\frac{3000}{600}=5\,\text{Вт}.
\]
Ответ: \(5\,\text{Вт}\).
}

% Source 687697: https://ege.fipi.ru/bank/questions.php?search=1&proj=BA1F39653304A5B041B656915DC36B38&qid=687697
\BankTask{687697}{
На сколько градусов нагреются \(100\,\text{г}\) свинца, если ему сообщить \(260\,\text{Дж}\) теплоты (в ${}^\circ\text{C}$)?
}{
\smallskip\noindent\textbf{Решение от GPT}\par
Используем табличную удельную теплоёмкость свинца \(c=130\,\text{Дж}/(\text{кг}\cdot\text{К})\). При \(m=0{,}1\,\text{кг}\)
\[
\Delta T=\frac{Q}{cm}=\frac{260}{130\cdot0{,}1}=20\,\text{К}.
\]
Ответ: \(20\,{}^\circ\text{C}\).
}

% Source 4EF342: https://ege.fipi.ru/bank/questions.php?search=1&proj=BA1F39653304A5B041B656915DC36B38&qid=4EF342
\BankTask{4EF342}{
Рабочее тело тепловой машины за цикл получает от нагревателя количество теплоты, равное \(100\,\text{Дж}\), и совершает работу \(60\,\text{Дж}\). Каков КПД тепловой машины (в $\%$)?
}{
\smallskip\noindent\textbf{Решение от GPT}\par
КПД тепловой машины:
\[
\eta=\frac{A}{Q_1}=\frac{60}{100}=0{,}6=60\%.
\]
Ответ: \(60\%\).
}

% Source 81C4BA: https://ege.fipi.ru/bank/questions.php?search=1&proj=BA1F39653304A5B041B656915DC36B38&qid=81C4BA
\BankTask{81C4BA}{
Рабочее тело тепловой машины с КПД \(30\%\) за цикл работы отдаёт холодильнику количество теплоты, равное \(70\,\text{Дж}\). Какое количество теплоты за цикл рабочее тело получает от нагревателя (в $\text{Дж}$)?
}{
\smallskip\noindent\textbf{Решение от GPT}\par
Из \(\eta=1-\dfrac{Q_2}{Q_1}\) получаем
\[
Q_1=\frac{Q_2}{1-\eta}=\frac{70}{1-0{,}30}=100\,\text{Дж}.
\]
Ответ: \(100\,\text{Дж}\).
}

% Source 0D600A: https://ege.fipi.ru/bank/questions.php?search=1&proj=BA1F39653304A5B041B656915DC36B38&qid=0D600A
\BankTask{0D600A}{
У идеального теплового двигателя Карно температура нагревателя равна \(227\,^{\circ}\mathrm{C}\), а температура холодильника равна \(27\,^{\circ}\mathrm{C}\). Определите КПД теплового двигателя (в $\%$).
}{
\smallskip\noindent\textbf{Решение от GPT}\par
Температуры резервуаров: \(T_1=227+273=500\,\text{К}\), \(T_2=27+273=300\,\text{К}\). КПД двигателя Карно:
\[
\eta=1-\frac{T_2}{T_1}=1-\frac{300}{500}=0{,}4=40\%.
\]
Ответ: \(40\%\).
}

% Source EC33AF: https://ege.fipi.ru/bank/questions.php?search=1&proj=BA1F39653304A5B041B656915DC36B38&qid=EC33AF
\BankTask{EC33AF}{
У идеального теплового двигателя Карно температура нагревателя \(400\,\text{К}\), температура холодильника \(300\,\text{К}\). За один цикл двигатель совершает работу \(16\,\text{кДж}\). Какое количество теплоты получает рабочее тело от нагревателя за один цикл (в $\text{кДж}$)?
}{
\smallskip\noindent\textbf{Решение от GPT}\par
КПД двигателя Карно:
\[
\eta=1-\frac{T_2}{T_1}=1-\frac{300}{400}=0{,}25.
\]
Из \(\eta=A/Q_1\) получаем:
\[
Q_1=\frac{A}{\eta}=\frac{16}{0{,}25}=64\,\text{кДж}.
\]
Ответ: \(64\,\text{кДж}\).
}

% Source 772CBF: https://ege.fipi.ru/bank/questions.php?search=1&proj=BA1F39653304A5B041B656915DC36B38&qid=772CBF
\BankTask{772CBF}{
Температура куска металла с удельной теплоёмкостью $900\,\text{Дж}/(\text{кг}\cdot\text{К})$ понизилась со $120\,{}^\circ\text{C}$ до $40\,{}^\circ\text{C}$. При этом выделилось количество теплоты, равное $108\,\text{кДж}$. Чему равна масса этого куска металла (в $\text{кг}$)?
}{
\smallskip\noindent\textbf{Решение от GPT}\par
$Q=cm(t_1-t_2)$, откуда $m=Q/[c(t_1-t_2)]=108000/(900\cdot80)=1{,}5\,\text{кг}$. Ответ: $1{,}5\,\text{кг}$.
}

% Source 8AB566: https://ege.fipi.ru/bank/questions.php?search=1&proj=BA1F39653304A5B041B656915DC36B38&qid=8AB566
\BankTask{8AB566}{
Какое количество теплоты необходимо для того, чтобы расплавить $2\,\text{кг}$ свинца, взятого при температуре плавления (в $\text{кДж}$)?
}{
\smallskip\noindent\textbf{Решение от GPT}\par
$Q=\lambda m=2{,}5\cdot10^4\cdot2=5\cdot10^4\,\text{Дж}=50\,\text{кДж}$. Ответ: $50\,\text{кДж}$.
}

% Source C2BD49: https://ege.fipi.ru/bank/questions.php?search=1&proj=BA1F39653304A5B041B656915DC36B38&qid=C2BD49
\BankTask{C2BD49}{
Тепловая машина с КПД $30\%$ за цикл работы получает от нагревателя количество теплоты, равное $120\,\text{Дж}$. Какую работу машина совершает за цикл (в $\text{Дж}$)?
}{
\smallskip\noindent\textbf{Решение от GPT}\par
$\eta=A/Q_1$, поэтому $A=\eta Q_1=0{,}30\cdot120=36\,\text{Дж}$. Ответ: $36\,\text{Дж}$.
}

% Source 80CCD9: https://ege.fipi.ru/bank/questions.php?search=1&proj=BA1F39653304A5B041B656915DC36B38&qid=80CCD9
\BankTask{80CCD9}{
Рабочее тело теплового двигателя за цикл совершает работу, равную $15\,\text{кДж}$, и получает от нагревателя количество теплоты, равное $75\,\text{кДж}$. Какое количество теплоты рабочее тело отдаёт холодильнику за цикл (в $\text{кДж}$)?
}{
\smallskip\noindent\textbf{Решение от GPT}\par
За цикл $\Delta U=0$, поэтому $Q_1=A+Q_2$. Отданная холодильнику теплота $Q_2=75-15=60\,\text{кДж}$. Ответ: $60\,\text{кДж}$.
}

% Source D1C429: https://ege.fipi.ru/bank/questions.php?search=1&proj=BA1F39653304A5B041B656915DC36B38&qid=D1C429
\BankTask{D1C429}{
Рабочее тело тепловой машины за один цикл отдаёт холодильнику количество теплоты $100\,\text{Дж}$ и совершает работу $30\,\text{Дж}$. Какое количество теплоты получает рабочее тело от нагревателя за два цикла (в $\text{Дж}$)?
}{
\smallskip\noindent\textbf{Решение от GPT}\par
За один цикл $Q_1=Q_2+A=100+30=130\,\text{Дж}$. За два цикла $Q=2Q_1=260\,\text{Дж}$. Ответ: $260\,\text{Дж}$.
}

% Source B44BA2: https://ege.fipi.ru/bank/questions.php?search=1&proj=BA1F39653304A5B041B656915DC36B38&qid=B44BA2
\BankTask{B44BA2}{
Тепловая машина с КПД $20\%$ совершает за цикл работу, равную $400\,\text{Дж}$. Какое количество теплоты машина получает за цикл от нагревателя (в $\text{Дж}$)?
}{
\smallskip\noindent\textbf{Решение от GPT}\par
$\eta=A/Q_1$, поэтому $Q_1=A/\eta=400/0{,}20=2000\,\text{Дж}$. Ответ: $2000\,\text{Дж}$.
}

% Source 0E8CD8: https://ege.fipi.ru/bank/questions.php?search=1&proj=BA1F39653304A5B041B656915DC36B38&qid=0E8CD8
\BankTask{0E8CD8}{
При сжатии идеального одноатомного газа при постоянном давлении внешние силы совершили работу $2000\,\text{Дж}$. Какое количество теплоты было передано при этом газом окружающим телам (в $\text{Дж}$)?
}{
\smallskip\noindent\textbf{Решение от GPT}\par
Работа газа $A=-2000\,\text{Дж}$. На изобаре $\Delta U=\frac32p\Delta V=\frac32 A=-3000\,\text{Дж}$. По первому началу $Q=\Delta U+A=-5000\,\text{Дж}$. Газ отдал $5000\,\text{Дж}$. Ответ: $5000\,\text{Дж}$.
}

\end{document}
